\chapter{Hillel Furstenberg \& Grigory Margulis (2020)}

\section{Biographical Sketches}



\subsection{Hillel Furstenberg}
Hillel Furstenberg was born on 29 September 1935 in Berlin, Germany. His family fled Nazi persecution and settled in the United States. He studied at Yeshiva University (BA, 1955) and Princeton University, where he completed his PhD in 1958 under the supervision of Salomon Bochner. After a brief appointment at the University of Minnesota, he joined the Hebrew University of Jerusalem in 1965, where he spent most of his career. Furstenberg received the Abel Prize in 2020, the Wolf Prize in 2007, and the Israel Prize in 1993.

\subsection{Grigory Aleksandrovich Margulis}
Grigory Margulis was born on 24 February 1946 in Moscow, Russia. He studied at Moscow State University, where he completed his PhD in 1970 under the supervision of Yakov Sinai. Despite facing significant restrictions as a Jewish mathematician in the Soviet Union, Margulis produced work of extraordinary depth. He held positions at Moscow State University and the Institute for Information Transmission Problems before emigrating to the United States in 1986, joining Yale University. Margulis received the Fields Medal in 1978, the Wolf Prize in 2005, and the Abel Prize in 2020.

\section{High-Level View for a General Audience}

\subsection{What Did Furstenberg and Margulis Do?}

Imagine you are studying the orbit of a point under a symmetry transformation---like following the path of a particle bouncing around inside a polygon. The overarching question: can we understand the long-term distribution of such orbits?

Furstenberg and Margulis developed the mathematical framework for studying such problems. They showed that for large classes of symmetric spaces, orbits of subgroups are either extremely regular (equidistributed) or extremely rare (discrete), with no intermediate possibilities.

\section{The Origin of the Problem}

The study of homogeneous spaces $G/\Gamma$ (where $G$ is a Lie group\index{Lie groups} and $\Gamma$ is a lattice) lies at the intersection of geometry, group theory, and dynamics.

\subsection{Early History}

The problem of understanding the distribution of orbits in homogeneous spaces dates back to the work of Gauss on binary quadratic forms and the reduction theory of lattices. In the 20th century, the study of flows on homogeneous spaces became central to number theory through the work of Hedlund, Artin, and Siegel.

\subsection{Furstenberg's Starting Point}

Furstenberg's work began in the 1960s with the study of distal flows---dynamical systems in which distinct points never approach each other under iterations. This led to a structure theory for minimal distal flows.

\subsection{Margulis' Starting Point}

Margulis began with the theory of discrete subgroups of Lie groups\index{Lie groups}, motivated by the work of Selberg, Borel, and Mostow on lattices in semisimple Lie groups\index{Lie groups}. The central questions were:
\begin{itemize}
\item Which groups admit lattices?
\item How are lattices classified?
\item What are the dynamics of group actions on homogeneous spaces?
\end{itemize}

\section{Deep Insight into the Problem}

\subsection{Furstenberg's Correspondence Principle}

Furstenberg realized that combinatorial problems about sets of integers can be transformed into problems about measure-preserving transformations. Specifically, given a set $A \subseteq \mathbb{Z}$ of positive density, we can construct:
\begin{itemize}
\item A probability space $(X,\mathcal{B},\mu)$.
\item An invertible measure-preserving transformation $T: X \to X$.
\item A measurable set $B \subseteq X$ such that $\mu(B) = d(A)$ and $n \in A \iff T^n x \in B$ for almost every $x$.
\end{itemize}

This transforms Szemer\'edi's theorem on arithmetic progressions into a theorem about multiple recurrence: for any set $B$ of positive measure, there exist $n \in \mathbb{N}$ such that $\mu(B \cap T^n B \cap T^{2n} B \cap \cdots \cap T^{(k-1)n} B) > 0$.

\subsection{Furstenberg's Boundary Theory}

Consider a random walk on a semisimple Lie group $G$ with step distribution $\mu$. The Furstenberg boundary is the maximal $G$-space on which the stationary measure is unique. For $G = SL(n,\mathbb{R})$, this is the flag manifold.

\subsection{Margulis' Measure Rigidity}

Margulis realized that invariant measures for group actions on homogeneous spaces must satisfy strong rigidity properties. A measure $\mu$ on $G/\Gamma$ that is invariant under a subgroup $H \subseteq G$ cannot be ``spread out'' unless $H$ is large enough.

\subsection{The Superrigidity Phenomenon}

Let $\Gamma$ be an irreducible lattice in $G = \prod_i G_i(k_i)$, a product of simple algebraic groups over local fields, with total rank at least 2. Margulis showed that any homomorphism $\Gamma \to H(k)$ (where $H$ is an algebraic group over a local field $k$) extends to a continuous homomorphism $G \to H(k)$, up to compact factors. This is called superrigidity.

\subsection{The Normal Subgroup Theorem}

For an irreducible lattice $\Gamma$ in a higher-rank semisimple Lie group $G$, any normal subgroup $N \trianglelefteq \Gamma$ is either:
\begin{itemize}
\item Contained in the center (hence finite), or
\item Of finite index in $\Gamma$.
\end{itemize}

This is remarkable: most groups have many normal subgroups (e.g., free groups have uncountably many), but higher-rank lattices are almost simple.

\subsection{The Oppenheim Conjecture}

Let $Q$ be an indefinite quadratic form in $n \geq 3$ variables that is not a scalar multiple of a rational form. The Oppenheim conjecture (1930) states that $Q(\mathbb{Z}^n)$ is dense in $\mathbb{R}$. Margulis proved this using unipotent flows on $SL(n,\mathbb{R})/SL(n,\mathbb{Z})$.

His insight: the values $Q(x)$ are related to the orbit of the identity under a one-parameter unipotent subgroup $u(t) \in SL(n,\mathbb{R})$:
\[
Q(x) = Q_0(g(t) \cdot \mathbb{Z}^n), \quad g(t) = \exp(tN),
\]
where $N$ is a nilpotent matrix chosen so that $u(t)$ preserves $Q_0$ but expands $\mathbb{Z}^n$. Then $Q(\mathbb{Z}^n)$ is the set of values of $Q_0$ at points in the orbit $u(t) \cdot \mathbb{Z}^n$. By Ratner's theorem, the closure of the orbit is itself a homogeneous space.

\section{Biggest Contributions}

\subsection{Furstenberg's Contributions}
\begin{itemize}
\item \textbf{Furstenberg's Correspondence Principle}: A deep connection between arithmetic progressions in combinatorics and measure-preserving transformations in ergodic theory\index{Ergodic theory}. This gave a completely new proof of Szemer\'edi's theorem and initiated the field of ergodic\index{Ergodic theory} Ramsey theory.
\item \textbf{Furstenberg's Boundary Theory}: The maximal boundary for random walks on Lie groups, classifying stationary measures. This became a fundamental tool in the study of random walks on groups and harmonic analysis on symmetric spaces.
\item \textbf{The Poisson Boundary}: A complete description of harmonic functions on symmetric spaces in terms of the Furstenberg boundary. This generalizes classical potential theory to non-amenable groups.
\item \textbf{Disjointness in Ergodic Theory\index{Ergodic theory}}: A concept analogous to disjointness in topological dynamics, classifying the structure of measure-preserving systems.
\item \textbf{Non-conventional Ergodic Averages}: The study of averages of the form:
\[
\frac{1}{N} \sum_{n=1}^N f_1(T^n x) f_2(T^{2n} x) \cdots f_k(T^{kn} x),
\]
leading to multiple recurrence theorems.
\end{itemize}

\subsection{Margulis' Contributions}
\begin{itemize}
\item \textbf{The Superrigidity Theorem}: Any irreducible lattice in a semisimple Lie group of rank at least 2 is arithmetic in the sense that it arises from the integer points of an algebraic group.
\item \textbf{The Normal Subgroup Theorem}: Normal subgroups of irreducible higher-rank lattices are either finite or of finite index, proving that these lattices are almost simple.
\item \textbf{The Oppenheim Conjecture}: For an indefinite quadratic form in $n \geq 3$ variables that is not proportional to a rational form, the set $Q(\mathbb{Z}^n)$ is dense in $\mathbb{R}$.
\item \textbf{The Margulis Lemma}: Any discrete subgroup of $SO(n,1)$ generated by sufficiently small rotations is nilpotent. This is a key ingredient in the study of hyperbolic manifolds.
\item \textbf{Margulis Numbers}: The Margulis constant for hyperbolic space: there exists $\varepsilon_n > 0$ such that in any hyperbolic $n$-manifold, every cyclic subgroup generated by a translation of length $<\varepsilon_n$ is central.
\item \textbf{Arithmeticity of Lattices}: Margulis completely classified the irreducible lattices in semisimple Lie groups, showing they are arithmetic when the group has rank at least 2.
\end{itemize}

\subsection{Unipotent Flows and Ratner's Theorems}

Building on Margulis' work on the Oppenheim conjecture, Marina Ratner proved a series of theorems that completely classify the invariant measures and orbit closures for unipotent flows on homogeneous spaces:
\begin{itemize}
\item \textbf{Measure Rigidity}: Any $H$-invariant and ergodic probability measure on $G/\Gamma$ (where $H$ is generated by unipotent one-parameter subgroups) is homogeneous: it is the $L$-invariant measure on a closed $L$-orbit for some closed subgroup $L$ containing $H$.
\item \textbf{Orbit Classification}: The closure of any $H$-orbit is itself a homogeneous space $Lx$ for some closed subgroup $L$.
\end{itemize}

\subsection{Ergodic Ramsey Theory}

Furstenberg's work spawned the field of ergodic Ramsey theory, which uses ergodic theory to prove combinatorial theorems:
\begin{itemize}
\item The Hales--Jewett theorem via ergodic theory.
\item The density Hales--Jewett theorem (Polymath project).
\item IP-sets and multiple recurrence.
\item Applications to van der Waerden's theorem and its generalizations.
\end{itemize}

\section{Biggest Contributions}

\subsection{Furstenberg's Contributions}
\begin{itemize}
\item \textbf{Furstenberg's Correspondence Principle}: A deep connection between arithmetic progressions in combinatorics and measure-preserving transformations in ergodic theory. This gave a completely new proof of Szemer\'edi's theorem.
\item \textbf{Furstenberg's Boundary Theory}: The maximal boundary for random walks on Lie groups, classifying stationary measures.
\item \textbf{The Poisson Boundary}: A complete description of harmonic functions on symmetric spaces.
\end{itemize}

\subsection{Margulis' Contributions}
\begin{itemize}
\item \textbf{The Superrigidity Theorem}: Any irreducible lattice in a semisimple Lie group of rank at least 2 is arithmetic (derived from the integer points of an algebraic group).
\item \textbf{The Normal Subgroup Theorem}: Normal subgroups of irreducible higher-rank lattices are either finite or of finite index.
\item \textbf{The Oppenheim Conjecture}: For an indefinite quadratic form in $n \geq 3$ variables, the set of values at integer points is either discrete or dense in $\mathbb{R}$.
\item \textbf{The Margulis Lemma}: Any discrete subgroup of $SO(n,1)$ generated by small rotations is nilpotent.
\end{itemize}

\section{Impact of the Discovery}

\subsection{Impact on Mathematics}

The work of Furstenberg and Margulis has transformed several areas of mathematics:

\begin{itemize}
\item \textbf{Number Theory}: The Oppenheim conjecture and its generalizations, the theory of quadratic forms, and the study of Diophantine approximation.
\item \textbf{Ergodic Theory}: The structure theory of measure-preserving systems, multiple recurrence, and ergodic Ramsey theory.
\item \textbf{Lie Theory}: The complete classification of lattices in semisimple Lie groups, the theory of discrete subgroups.
\item \textbf{Geometry}: The geometry of hyperbolic manifolds, the Margulis lemma, and the thick-thin decomposition.
\item \textbf{Combinatorics}: Furstenberg's correspondence principle revolutionized the study of arithmetic progressions and related structures.
\end{itemize}

\subsection{Impact on Other Fields}

The ergodic theory of group actions has found applications in:
\begin{itemize}
\item \textbf{Mathematical Physics}: The study of quantum chaos, the distribution of eigenfunctions, and quantum unique ergodicity.
\item \textbf{Computer Science}: Expander graphs and the theory of pseudo-randomness, building on the work of Margulis on expander families.
\item \textbf{Engineering}: The theory of random walks on groups and their applications to networks and algorithms.
\end{itemize}

\section{Current Developments}

\subsection{The Work of Eskin--Mirzakhani}

Alex Eskin and Maryam Mirzakhani used Ratner's theorems and Margulis' ideas to classify invariant measures on moduli spaces of flat surfaces.

\subsection{Arithmetic Quantum Unique Ergodicity}

Elon Lindenstrauss proved the arithmetic quantum unique ergodicity conjecture using tools from homogeneous dynamics.

\subsection{Approximate Groups}

The structure theory of approximate groups, developed by Breuillard, Green, and Tao, uses ideas from the theory of Lie groups and homogeneous dynamics.

\subsection{Superstrong Approximation}

Bourgain, Gamburd, and Sarnak developed superstrong approximation, using expansion properties of Cayley graphs of arithmetic groups.

\subsection{Applications to Diophantine Approximation}

Recent work by Kleinbock, Margulis, and others on Diophantine approximation on manifolds uses homogeneous dynamics.

\section{Conclusion}

Furstenberg and Margulis revolutionized homogeneous dynamics, creating deep connections between ergodic theory, group theory, number theory, and combinatorics. Their work continues to be extended and applied in diverse areas of mathematics.


\section{Behind the Breakthrough: The Human Story}
Margulis won the Fields Medal in 1978 but was denied a visa by the Soviet government to travel to Helsinki to receive it, causing an international outcry among mathematicians.

\section{Pedagogical Metaphors and Lineage}
\subsection{Signature Metaphor}
``Random walks on Lie group lattices that reveal rigid, symmetric properties of geometric spaces.''

\subsection{Mathematical Lineage}
Furstenberg advised by Salomon Bochner. Margulis advised by Pyotr Novikov.

\section{Applied Science and Computational Algorithms}
\subsection{Key Takeaways for Applied Science}
Expander graphs constructed using Margulis' techniques are used in building high-throughput routing networks, error-correcting codes, and hash functions.

\subsection{Algorithms and Numerical Implementations}
Belief Propagation decoding algorithms for Low-Density Parity-Check (LDPC) codes operate on expander graphs built from arithmetic group lattices.

\section{The Research Frontier}
Investigating random walks on general groups and lattices in Lie groups of rank 1, and establishing expander graph families from arithmetic subgroups.

\section{Guided Exercises and Challenges}
\begin{enumerate}[label=\textbf{\arabic*.}]
  \item Explain how Furstenberg gave an ergodic-theoretic proof of Szemer\'edi's theorem using the multiple recurrence theorem.
  \item State Margulis' superrigidity theorem for lattices in higher-rank semi-simple Lie groups.
  \item Explain the connection between expander graphs and Kazhdan's Property (T).
\end{enumerate}

\section{Curriculum Map and Further Reading}
For students wishing to delve deeper into the mathematical machinery underlying these breakthroughs, the following learning path is recommended:
\begin{itemize}
  \item H. Furstenberg, \emph{Recurrence in Ergodic Theory and Combinatorial Number Theory}, Princeton University Press.
  \item G. A. Margulis, \emph{Discrete Subgroups of Semisimple Lie Groups}, Springer.
\end{itemize}

\section*{Bibliography}
\begin{enumerate}
\item H. Furstenberg, "The structure of distal flows," American Journal of Mathematics 85 (1963), 477--515.
\item H. Furstenberg, "Recurrence in Ergodic Theory and Combinatorial Number Theory," Princeton University Press, 1981.
\item H. Furstenberg, "Nonconventional ergodic averages," in "The Legacy of Norbert Wiener," American Mathematical Society, 1997.
\item G. A. Margulis, "Discrete subgroups of semisimple Lie groups," Springer, 1991.
\item G. A. Margulis, "Oppenheim conjecture," in "Fields Medalists' Lectures," World Scientific, 1997.
\item M. Ratner, "On Raghunathan's measure conjecture," Annals of Mathematics 134 (1991), 545--607.
\item D. W. Morris, "Introduction to Arithmetic Groups," Deductive Press, 2015.
\item A. Eskin and M. Mirzakhani, "Invariant measures on the space of flat surfaces," arXiv:1302.1003, 2013.
\item E. Lindenstrauss, "Pointwise theorems for amenable groups," Inventiones Mathematicae 146 (2001), 259--295.
\item E. Breuillard, "A brief introduction to approximate groups," in "Thin Groups and Superstrong Approximation," Cambridge University Press, 2014.
\item Z. Rudnick and P. Sarnak, "The behaviour of eigenstates of arithmetic hyperbolic manifolds," Communications in Mathematical Physics 161 (1994), 195--213.
\item J. Bourgain and A. Gamburd, "Uniform expansion bounds for Cayley graphs of $SL_2(\mathbb{F}_p)$," Annals of Mathematics 167 (2008), 625--642.
\item E. Lindenstrauss, "Invariant measures and arithmetic quantum unique ergodicity," Annals of Mathematics 163 (2006), 165--219.
\item A. Eskin, "Quasi-isometric rigidity of nonuniform lattices in higher-rank symmetric spaces," Journal of the American Mathematical Society 11 (1998), 321--363.
\item G. A. Margulis, "Complete classification of discrete subgroups of semisimple Lie groups," in "Proceedings of the ICM 1978."
\end{enumerate}
